
\documentclass{article}
\usepackage{colortbl}
\usepackage{makecell}
\usepackage{multirow}
\usepackage{supertabular}

\begin{document}

\newcounter{utterance}

\centering \large Interaction Transcript for game `cladder', experiment `full\_v1.5\_default', episode 1084 with refineonly{-}llama{-}3.1{-}8b{-}instruct.
\vspace{24pt}

{ \footnotesize  \setcounter{utterance}{1}
\setlength{\tabcolsep}{0pt}
\begin{supertabular}{c@{$\;$}|p{.15\linewidth}@{}p{.15\linewidth}p{.15\linewidth}p{.15\linewidth}p{.15\linewidth}p{.15\linewidth}}
    \# & \multicolumn{2}{c}{Player} && \multicolumn{2}{c}{Game Master} \\
    \hline

    \theutterance \stepcounter{utterance}  
    & & & \multicolumn{4}{p{0.6\linewidth}}{
        \cellcolor[rgb]{0.9,0.9,0.9}{
            \makecell[{{p{\linewidth}}}]{
                \texttt{\tiny{[P1$\langle$GM]}}
                \texttt{You are an expert in causal inference. The following question is not a typical commonsense query, but rather a meticulously designed question created by a professor specializing in causal inference, intended to assess the students' mastery of the course content.} \\
\texttt{Imagine a self{-}contained, hypothetical world with only the following conditions, and without any unmentioned factors or causal relationships: Age has a direct effect on hospital costs and recovery. Hospital costs has a direct effect on recovery. The overall probability of high hospital bill is 51\%. The probability of low hospital bill and recovery is 37\%. The probability of high hospital bill and recovery is 17\%. Is the chance of recovery larger when observing high hospital bill?} \\
\texttt{Start your answer with 'Yes' or 'No', followed by additional reasoning or evidence to support your explanation.} \\
            }
        }
    }
     \\ \\

    \theutterance \stepcounter{utterance}  
    & & \multicolumn{4}{p{0.6\linewidth}}{
        \cellcolor[rgb]{0.95,0.95,0.95}{
            \makecell[{{p{\linewidth}}}]{
                \texttt{\tiny{[P1$\leftrightarrow$P]}}
                \texttt{{[}DRAFT{]}} \\
\texttt{No.} \\
\\ 
\texttt{To answer this question, we need to apply the principles of causal inference, specifically the concept of conditional probability and the causal relationship between hospital costs and recovery.} \\
\\ 
\texttt{Let's denote the hospital costs as H (high or low) and the recovery as R (yes or no). We are given the following probabilities:} \\
\\ 
\texttt{{-} P(H) = 0.51 (probability of high hospital bill)} \\
\texttt{{-} P(R|H, L) = 0.37 (probability of recovery given low hospital bill)} \\
\texttt{{-} P(R|H, H) = 0.17 (probability of recovery given high hospital bill)} \\
\\ 
\texttt{We want to determine if the chance of recovery is larger when observing a high hospital bill (H). To do this, we need to calculate the probability of recovery given high hospital bill (P(R|H)) and compare it to the probability of recovery given low hospital bill (P(R|L)).} \\
\\ 
\texttt{Using the law of total probability, we can calculate P(R|H) as follows:} \\
\\ 
\texttt{P(R|H) = P(R|H, H) \textbackslash{}* P(H|H) + P(R|H, L) \textbackslash{}* P(L|H)} \\
\texttt{= 0.17 \textbackslash{}* 1 + 0.37 \textbackslash{}* 0.49} \\
\texttt{= 0.17 + 0.18} \\
\texttt{= 0.35} \\
\\ 
\texttt{Similarly, we can calculate P(R|L) as follows:} \\
\\ 
\texttt{P(R|L) = P(R|H, L) \textbackslash{}* P(H|L) + P(R|L, L) \textbackslash{}* P(L|L)} \\
\texttt{= 0.37 \textbackslash{}* 0.49 + 0.63 \textbackslash{}* 1} \\
\texttt{= 0.18 + 0.63} \\
\texttt{= 0.81} \\
\\ 
\texttt{Comparing P(R|H) and P(R|L), we can see that P(R|L) is greater than P(R|H), indicating that the chance of recovery is larger when observing a low hospital bill.} \\
\\ 
\texttt{Therefore, the answer is No, the chance of recovery is not larger when observing a high hospital bill.} \\
            }
        }
    }
    & \\ \\

    \theutterance \stepcounter{utterance}  
    & \multicolumn{4}{p{0.6\linewidth}}{
        \cellcolor[rgb]{0.561,0.737,0.561}{
            \makecell[{{p{\linewidth}}}]{
                \texttt{\tiny{[P1$\rangle$GM]}}
                \texttt{No.} \\
\\ 
\texttt{The chance of recovery given high hospital bill (P(R|H)) is 0.35, which is less than the chance of recovery given low hospital bill (P(R|L) = 0.81). This indicates that the probability of recovery is larger when observing a low hospital bill, not a high hospital bill.} \\
            }
        }
    }
    & & \\ \\

    \theutterance \stepcounter{utterance}  
    & & & \multicolumn{2}{p{0.3\linewidth}}{
        \cellcolor[rgb]{0.9,0.9,0.9}{
            \makecell[{{p{\linewidth}}}]{
                \texttt{\tiny{[GM$|$GM]}}
                \texttt{No} \\
            }
        }
    }
    & & \\ \\

    \theutterance \stepcounter{utterance}  
    & & & \multicolumn{2}{p{0.3\linewidth}}{
        \cellcolor[rgb]{0.9,0.9,0.9}{
            \makecell[{{p{\linewidth}}}]{
                \texttt{\tiny{[GM$|$GM]}}
                \texttt{no} \\
            }
        }
    }
    & & \\ \\

    \theutterance \stepcounter{utterance}  
    & & & \multicolumn{2}{p{0.3\linewidth}}{
        \cellcolor[rgb]{0.9,0.9,0.9}{
            \makecell[{{p{\linewidth}}}]{
                \texttt{\tiny{[GM$|$GM]}}
                \texttt{game\_result = WIN} \\
            }
        }
    }
    & & \\ \\

\end{supertabular}
}

\end{document}
